\documentclass[10pt]{article}
\usepackage[margin=20mm]{geometry}
\usepackage{amsmath,amssymb,mathtools,fontspec,unicode-math,xurl,hyperref}
\IfFontExistsTF{DejaVu Serif}{\setmainfont{DejaVu Serif}}{\setmainfont{Latin Modern Roman}}
\setmathfont{latinmodern-math.otf}
\setlength{\parindent}{0pt}\setlength{\parskip}{6pt}\sloppy
\title{A technical review of
The multiplier spectrum morphism is generically injective}
\author{AI-assisted technical review}
\date{27 October 2026}
\begin{document}
\maketitle
Abstract: This review examines the exact preprint arXiv:2309.15382v3. The definition of intertwining permits a union of a horizontal and a vertical invariant line. Consequently the nonempty generic rigidity assertion used in the main proof cannot hold under that definition. A separate problem concerns the all-critical-point PCE and PR assertions in the dynamical André–Oort argument: persistently periodic critical sections violate them identically. These are Major issues in the stated proof chain, with specific, bounded repair routes. A missing cycle-length factor in the elementary counting lemma is a Minor error whose correction preserves its conclusion. None of these findings gives a counterexample to generic injectivity of the multiplier spectrum. Independent scientific adjudication remains pending.\par
\section*{Identity, scope, and proof structure}
The reviewed work is The multiplier spectrum morphism is generically injective, by Zhuchao Ji and Junyi Xie [target], frozen identity M25-0366. The exact source has 22 PDF pages and is arXiv:2309.15382v3, submitted on 22 September 2025; its first page separately prints 23 September 2025. The associated DOI is https://doi.org/10.4171/jems/1750. The journal record dates online publication to 7 November 2025. No equality between the journal full text and the reviewed preprint is presumed. Page numbers below are one-based PDF pages, agreeing with the printed numbers.\par
All 22 pages, including proofs, the alternative approach, and references, were read as extracted text. Twelve source pages and six reference pages were also actually inspected as images. This is an initial scientific review, neither an independent acceptance nor a human review.\par
For degree \(d\geq2\), \(S_n\) encodes the multipliers of all \(d^n+1\) fixed points of \(f^n\), counted with fixed-point multiplicity. The descending closed equivalence relations defined by the truncated data stabilize by Noetherianity. The resulting \(\tau_d\) must have a singleton full fiber at a generic point (Definition 1.2). The theorem is generic injectivity, not injectivity everywhere.\par
For \(d\geq4\), the proof combines disjoint-type PCF parameters, curves with \(2d-3\) prescribed superattracting cycles, a non-diagonal finite étale fiber relation for \(\tau_d\), a two-family DAO argument, and generic rigidity of intertwined simple maps. Degrees two and three are treated by external results. The stated conjectures about small periods, length spectra, and the full exceptional locus are kept separate from the proved claims.\par
\section*{Substantive findings}
Problem 1 (The reducible-curve definition makes the rigidity premise universal).  Definition 3.1, page 11, permits a reducible invariant curve whose two projections are surjective. Under this definition every pair of rational maps is intertwined. Thus the nonempty form of Theorem 3.3 used on pages 15 and 17 is false.\par
Complete counterconstruction. Let \(d\geq4\), and let \(f,g\in\operatorname{Rat}_d(\mathbb C)\) be arbitrary. Each map has a fixed point on \(\mathbb P^1\); choose \(f(p)=p\) and \(g(q)=q\). Set \[Z=(\{p\}\times\mathbb P^1)\ \cup\ (\mathbb P^1\times\{q\}).\] This is a closed, reduced, pure one-dimensional algebraic curve. Both projections of the whole \(Z\) are surjective. Since a nonconstant rational self-map of \(\mathbb P^1\) is surjective, \((f\times g)(Z)=Z\). Thus every requirement of Definition 3.1 holds, including equality invariance.\par
Suppose the open set \(W_d\) used in Theorem 3.3 were nonempty, and choose \(f\in W_d\). Choose \(g\) in the simple-map locus outside the closure of the conjugacy orbit of \(f\). Such a \(g\) exists: the simple locus is a nonempty open subset of the \(2d+1\)-dimensional variety \(\operatorname{Rat}_d\), whereas a \(\operatorname{PGL}_2\)-orbit and its closure have dimension at most three. Now \(f\in W_d\), \(g\) is simple, and \(f,g\) are intertwined by the displayed \(Z\), but they are not conjugate. This construction works for every possible \(f\) in any proposed nonempty \(W_d\); it does not assume that a chosen numerical map belongs to an unspecified generic set.\par
Impact and alternative readings. Theorem 3.3 does not explicitly say that \(W_d\) is nonempty, so the empty set would make its isolated existential statement vacuous. That reading cannot supply the explicitly nonempty open \(U\) on page 15. Reading surjectivity component by component adds a condition absent from Definition 3.1. The issue affects the central rigidity implication used to close the main contradiction. It is Major, even though the necessary correction is short.\par
Bounded repair. Require finite projections, or a periodic irreducible component dominating both factors. A finite product self-map permutes the components of an invariant finite union, so a component is invariant under an iterate. The consulted version of Pakovich’s Theorem 1.4 [pak], page 3, expressly assumes an irreducible curve that is neither horizontal nor vertical. It applies to this corrected situation, followed by the same-entropy rigidity step.\par
The objects \(\operatorname{Corr}(\mathbb P^1)^{f_t\times g_t}_*\) actually introduced on page 14 have finite projections. Thus the DAO construction aims at a stronger object capable of supporting the repair. The correction changes the definition and the valid theorem interface; it does not prove the original broad assertion. No counterexample to Theorem 1.3 is established.\par
Problem 2 (The all-critical-point typicality assertions are false).  The PCE and PR definitions on page 12 quantify over every marked critical point. Assertions (4) and (5) on page 14 then claim these properties almost everywhere under Theorem 3.4’s hypotheses. These claims fail for admissible families and for the periodic critical sections in the main construction.\par
An explicit family satisfying all theorem hypotheses. Take \(V=\mathbb A^1_\mathbb C\) and \(f_t(z)=g_t(z)=z^2+t\). These are algebraic degree-two families, non-isotrivial, and not flexible Lattès families. The product of the two finite fixed-point multipliers is \(4t\), which verifies non-isotriviality. Equality of the two families is allowed.\par
There are infinitely many simultaneous PCF parameters. Put \[P_1(t)=t,\qquad P_{n+1}(t)=P_n(t)^2+t.\] Then \(P_n(t)=f_t^n(0)\), its degree is \(2^{n-1}\), and its linear coefficient is one. For each prime \(p\), \(P_p\) has a simple zero at zero and at least one other complex zero. At a nonzero zero, \(0\) has exact period \(p\): its period divides \(p\) and cannot be one. Different primes give distinct parameters. Infinity is fixed, so these parameters are PCF.\par
In the coordinate \(w=1/z\), the family near infinity is \[w\longmapsto \frac{w^2}{1+tw^2}.\] The marked critical section at infinity is constant and fixed. At every iterate its parameter derivative and its distance to the critical set are zero. Neither a positive exponential PCE lower bound nor the PR lower bound can hold at any parameter. The obstruction is not a null exceptional set.\par
Even restricting to the active section \(c_0(t)=0\) leaves the printed PCE condition at \(n=0\) false: \[\frac{d}{dt}f_t^0(c_0(t))=0.\] A finite-prefix convention is necessary as well as the restriction on critical points.\par
Direct relevance to the main families. For a persistently periodic critical section \(c\) of period \(r\), \[\xi_{c,kr}(t)=c(t),\qquad
 \left\|d\xi_{c,kr}(t)/dt\right\|=\left\|dc(t)/dt\right\|.\] At a fixed parameter the right-hand side is finite and independent of \(k\); it cannot dominate \(C\lambda^{kr}\) with \(C>0,\lambda>1\). At the same iterates the distance to the critical set is zero. The curves in Lemma 2.5 and (3.1) have persistently periodic critical points after marking. Thus the literal conditions (4),(5) are incompatible with the families used by the main argument.\par
Dependency check and repair boundary. The consulted DAO preprint [dao] defines conditions for a marked pair. Its proof on page 52, condition (b), restricts them to active critical points; Proposition 7.8, page 50, uses the selected marked points. This supports a repair using only the active \(a_1,b_1\). The cited preprint’s Definition 4.1 also prints \(n\geq0\), so citing that definition alone does not fix the constant-section obstruction. One must allow a finite initial prefix, or pass to iterated marked sections and check measure scaling and renormalization hypotheses. The last paragraph on target page 14 also switches from \(b_1\) to \(b_2\); the active point must be kept consistent.\par
This is a Major issue in the stated DAO proof: the purported full-measure set supplying its inputs can be empty. The contradictions are established; a complete replacement of the deep two-family argument is not claimed. The diagonal intertwines the explicit example even in the stronger sense, so it is not a counterexample to Theorem 3.4 or generic injectivity. Section 3.4 is an outline, not an independently expanded replacement proof in this review.\par
Problem 3 (The cycle-count inversion omits cycle length).  The two counting formulas at the end of Lemma 3.5’s proof, page 16, confuse cycles with fixed points.\par
Write \(u_n\) for the number of zero multipliers among fixed points of \(f^n\), and \(m_\ell\) for the number of critical cycles of exact length \(\ell\). Each such cycle contributes \(\ell\) points when \(\ell\mid n\). These are simple fixed points of \(f^n\), since a zero multiplier makes the derivative of the local fixed-point equation equal to \(-1\). Thus \[u_n=\sum_{\ell\mid n}\ell m_\ell,\qquad
 m_n=\frac1n\sum_{\ell\mid n}\mu(\ell)u_{n/\ell}.\] The printed formulas omit \(\ell\) and \(1/n\).\par
For \(f(z)=z^2-1\), the critical cycles are fixed infinity and \(0\mapsto-1\mapsto0\). This is a degree-two hyperbolic PCF map of disjoint type \((1,2)\); choosing \(g=f\) also meets the same-spectrum premise. Here \(u_1=1,u_2=3\), while \(m_2=1\). The printed inversion gives two. An independent exact polynomial calculation of the greatest common divisors of \(f^n(z)-z\) and \((f^n)'(z)\), with infinity added, verified \[(u_1,u_2,u_3,u_4)=(1,3,1,3),\qquad
 (m_1,m_2,m_3,m_4)=(1,1,0,0).\] This is Minor: the corrected identities still determine every critical cycle length. Having \(2d-2\) distinct superattracting cycles exhausts the critical-point budget, so Lemma 3.5’s intended conclusion survives.\par
\section*{Remaining coverage and qualifications}
The Noetherian truncation and full-fiber injectivity were checked. In Lemma 2.5, application of the curve lemma to a higher-dimensional critical-relation locus requires a non-isotrivial curve slice. If the remaining critical point were persistently preperiodic, a forbidden non-isotrivial PCF family would already result. These are relevant implicit steps; no counterexample to that lemma was constructed.\par
The off-diagonal finite étale relation in Section 3.3 was followed through its curve lifts and common base change. Dominant curve maps may be restricted to dense opens. The equality of spectra and inequality of conjugacy classes were distinguished from the faulty final rigidity interface. The reduction to \(\overline{\mathbb Q}\), active bifurcation measures, transcendental-parameter lifting, and CE/TCE implications were read with their external dependencies retained.\par
The polynomial proof is brief. Reusing the rational-map count of prescribed relations requires adaptation to the number of finite polynomial critical points. Its cited rigidity results and the full Favre–Gauthier monograph were not independently reconstructed. This review also does not independently reprove the full arithmetic equidistribution, every distortion estimate in the DAO paper, the low-degree external results, or Remark 1.6’s Berkovich specialization sketch. These are coverage limits, not asserted counterexamples.\par
A bounded search of the exact arXiv and journal records and of the title with code, supplement, and correction terms found no applicable supplementary artifact or erratum. Absence and priority are not established. No author code was executed. The two reference PDFs were consulted only in the recorded passages; later published versions are not silently substituted.\par
The supported overall assessment is Major, arising from Problems 1 and 2. Problem 3 has a complete elementary correction. This review has not disproved the main generic-injectivity theorem.\par
\section*{Supporting checks}
The bound source is `papers/0366/0366.pdf`, 456,735 bytes, SHA256\par
`68e963013001203ef3d365c9e5c2824f0141b7a7d9ddcdf74492e2aecb5356c6`.\par
Reviewer-authored SymPy 1.13.3 calculations checked the cycle-count example, the infinity coordinate map, and finite instances of the PCF parameter recursion. The infinitely-many-prime-periods argument is analytic, not an inference from finitely many computations.\par
The report uses the existing official Tectonic 0.17.0 runtime with shell escape disabled. The delivery record binds exact TeX, compiler, dependency payloads before and after compilation, final PDF, parsing, and actual visual inspection of every final report page. Rendering alone is not inspection. Earlier incomplete states and revisions remain in the history.\par
Scientific reasoning, drafting, and visual inspection belong to the current scientific conversation. Native tools perform the recorded byte operations, calculations, compilation, and rendering. Main’s role is UI/logistics and byte transport. The controller-observed Academic Researchers / Chat / 6 Pro label is separate from the runtime self-description GPT-6 Astra Pro; backend equivalence was not independently verified. Scientific decision is pending; human reviews are empty.\par
9 Z. Ji and J. Xie, The multiplier spectrum morphism is generically injective, arXiv:2309.15382v3, 22 September 2025, 22 pages. https://arxiv.org/abs/2309.15382v3. F. Pakovich, On iterates of rational functions with maximal number of critical values, arXiv:2107.05963v2, 10 January 2023. Consulted PDF pages 3 and 15; Theorem 1.4 visually checked. https://www.math.bgu.ac.il/\textasciitilde{}pakovich/Publications/imax.pdf. Z. Ji and J. Xie, DAO for curves, arXiv:2302.02583v2, 27 September 2023. Consulted PDF pages 22–25 and 50–53. https://arxiv.org/abs/2302.02583v2.\par
\end{document}
